\documentclass[12pt]{article}
\usepackage[utf8]{inputenc}
\usepackage[T1]{fontenc}
\usepackage{hyperref}
\title{Finongan Language Resources for K-12 Curriculum Development}
\author{Compiled by Claude Sonnet 5.5 (claude-sonnet-5-5) for Andr\'as Kornai}
\date{}
\begin{document}
\maketitle

\section{Introduction}
Finongan (Finungwa; ISO 639-3: fag; Glottocode fino1238) is a Finisterre-Huon language of Morobe Province, Papua New Guinea. Speaker counts conflict: 1,140 (2006), about 1,200 in two villages of the Lowai area (1988), about 1,500 (dictionary project), 2,500--2,600 (Joshua Project), 470 (Global Recordings Network). Wikitongues rates it vigorous (EGIDS 6a). A dictionary exists on Webonary, and Chris and Amy Rice have collected data since 2000.

\section{Primary Educational Resources}
No primers, textbooks or school programmes were found in search results.

\section{Linguistic and Reference Materials}
\begin{itemize}
\item Finongan dictionary on Webonary, with a foreword describing two dialects. \url{https://www.webonary.work/finongan/overview/foreword/}
\item Wikipedia, Finongan language. \url{https://en.wikipedia.org/wiki/Finongan_language}
\item Counting-systems study page (MPI EVA channumerals). \url{https://lingweb.eva.mpg.de/channumerals/Finongan.htm}
\end{itemize}

\section{Religious and Community Resources}
\begin{itemize}
\item Joshua Project profile of Finungwan; about ninety percent Christian, with a Bible portion translation reported. \url{https://joshuaproject.net/people_groups/11210}
\end{itemize}

\section{Digital and Audio Resources}
\begin{itemize}
\item Global Recordings Network, Finongan. \url{https://globalrecordings.net/en/language/fag}
\item Wikitongues, Finongan. \url{https://wikitongues.org/languages/fag/}
\end{itemize}

\section{Curriculum Development Context}
The language is stable and used by all in the community, but is not known to be taught in schools. The Webonary dictionary is the main basis for a native-language curriculum; the body-tally numeral system is relevant to mathematics teaching. Speaker figures should be checked locally.

\end{document}